\section{Research Objectives}

This thesis studies lifelong model editing under a sequential update regime. The first objective is to examine how representative editing methods behave when new facts arrive incrementally rather than being optimized together in a single batch. The second objective is to characterize the trade-off between editing quality and operational cost, especially the gap between methods that achieve strong factual rewriting and methods that are easier to imagine in a continuously updated environment.

To support that objective, the empirical study evaluates editing behavior along three dimensions: whether the revised fact is expressed reliably across prompt variants, whether unrelated knowledge is preserved, and whether the revised fact remains usable in downstream compositional questions. These dimensions are formalized in Chapter 4 and serve as the basis for comparing the selected baselines.
